The electrification of heating and cooling is connecting more thermostatically controlled loads (TCLs), such as heat pumps (HPs), to low-voltage networks. Distribution system operators (DSOs) need the installed TCL capacity behind each feeder to plan for cold-weather peaks, but rarely know it. This paper estimates this capacity from the feeder net load and outdoor temperature a DSO already holds, plus the temperature sensitivity of the simultaneity factor (SF). The SF is the share of installed capacity in use, and its sensitivity is how fast this share grows as the temperature falls. The SF sensitivity is calibrated once on a small pilot of households with submetered heating. Dividing the feeder's net-load sensitivity, the rate at which its net load grows as the temperature falls, by the SF sensitivity gives its installed capacity. On 1260 semi-synthetic aggregates of Swiss households, the capacity is recovered at a weighted absolute percentage error (WAPE) of 19\%, with the error growing as the share of households with an HP falls. On three real Swiss and German aggregates outside the pilot, where every household has an HP, the estimates lie within 10\% of the true capacity. The same data recover the temperature-driven HP energy, with no pilot, at a WAPE of 9\%, and a flexibility potential at a WAPE of 10\%. Across populations from five countries, the SF sensitivity transfers within 12\% inside one climate zone. It differs by 13--14\% between HP technologies, however, and up to 1.7-fold between climate zones and between social and private housing.
